\documentclass[11pt]{article}
\usepackage{amsmath,amssymb,amsthm}
\usepackage[margin=1in]{geometry}
\title{Nonexistence of integer solutions to $6+x^{3}-y^{2}+y^{2}z^{2}=0$}
\author{}
\date{}

\begin{document}
\maketitle

\section*{Theorem}
There are no integers $x,y,z$ with $z$ even, $y$ odd, and $x\equiv 3\pmod 4$ satisfying
\[6 + x^{3} - y^{2} + y^{2}z^{2}=0.\]

\section*{Proof}
Rewrite the equation as
\begin{equation}\label{eq:main}
x^{3}+6 = -y^{2}(z^{2}-1).
\end{equation}

Set $X := -x$. Since the right-hand side of \eqref{eq:main} is negative when
$y\neq 0$ and $|z|\ge 2$, any solution with $y\neq 0$ must satisfy
$x^{3}+6<0$. Thus $x\le -3$, and because $x\equiv 3\pmod 4$ we have
$x\le -5$ and hence $X\ge 5$ with $X\equiv 1\pmod 4$.

Rewrite \eqref{eq:main} as
\begin{equation}\label{eq:Xeq}
X^{3}+6 = y^{2}(z^{2}-1),
\end{equation}
where $X\in\mathbb Z_{>0}$, $X\equiv 1\pmod 4$, $y$ is odd, and $z$ is even.

\paragraph{1) Coprimality of $X$ and $y$.}
Suppose a prime $p$ divides both $X$ and $y$. Since $y$ is odd, $p$ is odd.
Then $p$ divides $X^{3}+6$ and therefore divides $6$. Thus $p=3$.
Write $X = 3X'$ and $y = 3y'$. Substituting in \eqref{eq:Xeq} gives
\[27X'^{3}+6 = 9y'^{2}(z^{2}-1).
\]
Dividing by $3$ yields
\[9X'^{3}+2 = 3y'^{2}(z^{2}-1).
\]
The left-hand side is congruent to $2$ modulo $3$, while the right-hand side is
congruent to $0$ modulo $3$. This is impossible. Hence no prime divides both
$X$ and $y$, i.e.
\begin{equation}\label{eq:coprime}
\gcd(X,y)=1.
\end{equation}

\paragraph{2) Parity and modulo~4 properties.}
Because $z$ is even, both $z-1$ and $z+1$ are odd. Their product
$z^{2}-1$ is therefore odd, and $y^{2}(z^{2}-1)$ is odd. In particular,
$X^{3}+6$ is odd, so $X$ is odd. Combined with $X\equiv 1\pmod 4$, we obtain

\begin{equation}\label{eq:mod4}
X+2 \equiv 3 \pmod 4,
\qquad
X^{2}-2X+4 \equiv 3 \pmod 4.
\end{equation}

The factorization
\begin{equation}\label{eq:factor}
X^{3}+6 = (X+2)(X^{2}-2X+4)
\end{equation}
takes the form of two coprime odd factors. Indeed,
\begin{align*}
\gcd(X+2, X^{2}-2X+4)
&= \gcd(X+2, (X^{2}-2X+4) - (X-4)(X+2)) \\
&= \gcd(X+2, 12).
\end{align*}
Because $X\equiv 1\pmod 4$, the number $X+2$ is an odd integer congruent to
$3$ modulo $4$. Thus $\gcd(X+2,12)$ divides $3$.

If $3 \mid X+2$, then $X\equiv 1\pmod 3$ and hence $X^{3}+6\equiv 1\pmod 3$.
But $z$ is even, so $z\equiv 0$ or $2 \pmod 3$, and therefore
$z^{2}-1 \equiv 2$ or $0 \pmod 3$. Since $y$ is odd, $y^{2}\equiv 1\pmod 3$.
Thus the right-hand side of \eqref{eq:Xeq} is congruent to $2$ or $0$ modulo $3$,
contradicting $X^{3}+6\equiv 1\pmod 3$. Consequently $3\nmid X+2$ and we have

\begin{equation}\label{eq:gcd1}
\gcd(X+2, X^{2}-2X+4) = 1.
\end{equation}

\paragraph{3) The parity of $z\pm 1$.}
Since $z$ is even, $z-1$ and $z+1$ are odd and coprime. One of them is
congruent to $1$ modulo $4$ and the other to $3$ modulo $4$.
(Indeed, any even $z$ satisfies $z\equiv 0$ or $2 \pmod 4$, so
$z-1 \equiv 3$ and $z+1 \equiv 1$, or vice versa.)

Because $y$ is odd, $y^{2}\equiv 1 \pmod 4$.
Therefore any odd integer obtained by multiplying one of $z-1$ or $z+1$ by
$y^{2}$ has the same residue class modulo $4$ as the original factor.
In particular,
\begin{equation}\label{eq:rhs-mod4}
\{(z-1)y^{2}, (z+1)y^{2}\} \equiv \{1,3\} \pmod 4.
\end{equation}

\paragraph{4) Final contradiction.}
Equation~\eqref{eq:Xeq} can be rewritten using \eqref{eq:factor} as
\begin{equation}\label{eq:split}
(X+2)(X^{2}-2X+4) = y^{2}(z-1)(z+1).
\end{equation}
The left-hand side is a product of two coprime odd integers by
\eqref{eq:gcd1}. Since $z-1$ and $z+1$ are coprime, the prime divisors of
$(z-1)(z+1)$ must split between the two coprime left factors without overlap.
Because $y^{2}$ is a square and every odd square is congruent to $1$ modulo $4$,
the residue of each factor modulo $4$ is determined entirely by the odd part
coming from either $z-1$ or $z+1$.

One of $z-1$ and $z+1$ is $1\pmod 4$ and the other is $3\pmod 4$. Hence in
any factorisation of $y^{2}(z-1)(z+1)$ into two coprime odd factors, one of the
resulting factors must be congruent to $1$ modulo $4$ and the other to $3$ modulo
$4$. This follows from the fact that an odd square is $1$ modulo $4$, so the
parity class of each factor is unchanged by the square multiplier.

However, from \eqref{eq:mod4} both factors on the left-hand side of
\eqref{eq:split} are congruent to $3$ modulo $4$. This is impossible if the
right-hand side splits as one factor congruent to $1$ and one factor congruent
to $3$ modulo $4$. Therefore no integer solution can exist.

\paragraph{5) Conclusion.}
The assumptions $z$ even, $y$ odd, and $x\equiv 3 \pmod 4$ force
$x\le -5$, $X=-x\ge 5$, $X\equiv 1\pmod 4$, and coprimality
$\gcd(X,y)=1$. The factorisation argument above then yields a rigid mod-$4$
contradiction. Hence the original equation has no integer solution under the
given hypotheses.

\end{document}
